\section{The Dynamics of Persistence}
\label{sec:dynamics_of_persistence}

To describe how a system maintains its architecture through time, we first identify the operational requirements of its state transitions: the work performed per update, the cost of that work, and the frequency of updates. These three quantities define a local operational cost rate. We then examine the constraints imposed by a finite resource budget to measure the cumulative cost of a history.

\subsection{State Transition Costs}
\label{subsec:state_transition_costs}

The operation of the six pillars through time introduces the constraint that transitions must be realized within the substrate's available resources and response times. Quantifying this requires distinguishing the work an update requires, the cost of that work, and the update frequency.

\begin{enumerate}
    \item \textbf{Structural workload ($K$):}
    The number of elementary operations required for one maintenance update of the local structure. An update preserves, repairs, or modifies the organization through which the pillars are instantiated. Here, $K$ measures operations per update. Its relationship to algorithmic description length \cite{kolmogorov_three_1968} is implementation-specific.
    \textit{Domain Manifestation (Computational Workload):} In a computational implementation, $K$ counts the elementary operations executed by an update procedure under a specified computational model.

    \item \textbf{Operational toll ($c$):}
    The mean resource cost per elementary operation, expressed in a specified accounting unit. Substrate manifestations include dissipated energy, computational expenditure, or other explicitly defined resources. Transition latency imposes a separate constraint on achievable update rates; it is not automatically equivalent to dissipative cost.
    \textit{Domain Manifestation (Information Erasure):} In the standard isothermal setting, erasing one unbiased bit requires at least $k_B T\ln 2$ of heat dissipation \cite{landauer_irreversibility_1961}. This provides a physical lower bound on the toll for that class of operation.

    \item \textbf{Coordinated transition flux ($f$):}
    The number of maintenance updates performed per unit substrate measure per unit time. An update involves coordinated operations, whose total is recorded by $K$. The flux required for persistence is dictated by the system's organization and the environmental disturbances it must accommodate.
    \textit{Domain Manifestation (Communication-Constrained Control):} Data-rate theorems establish communication requirements for stabilizing specified dynamical systems \cite{nair_feedback_2007}. For updates carrying a fixed amount of transmitted information, such requirements can constrain update frequency.
\end{enumerate}

While a system might be made up of sub systems each performing transitions, the parent system transition costs will only be the sum of the subsystem costs.


\subsection{Operational Cost Rate Density}
\label{subsec:dissipation_density}

Repeated operations draw on a finite resource budget. Comparing this demand with available resources requires combining the workload, operational toll, and update flux into a local cost rate per unit substrate measure.

\begin{align*}
[c] &= \frac{\text{cost}}{\text{operation}} \\[6pt]
[K] &= \frac{\text{operations}}{\text{update}} \\[6pt]
[f] &= \frac{\text{updates}}{\text{substrate measure} \cdot \text{time}} \\[6pt]
\end{align*}

For transitions requiring $K$ elementary operations each, if the mean cost of those operations is $c$, the cost per update is $K c$. Performing updates at flux $f$ therefore yields the local density:
\begin{equation}
\label{eq:dissipation_density}
    \mathcal{D}=cKf,
    \qquad
    [\mathcal{D}]
    =\frac{\text{cost}}
    {\text{substrate measure}\cdot\text{time}}
\end{equation}

The interpretation of $\mathcal{D}$ depends on the resource used to define $c$. When $c$ measures dissipated energy per operation, $\mathcal{D}$ is a dissipation density; otherwise it records the
corresponding operational resource cost.

When any operational cost is 0, $\mathcal{D}=0$. Whether the corresponding inactivity or absence of expenditure compromises persistence depends on the maintenance required over the observation interval.





\subsection{The Complexity Trade-off}
\label{subsec:complexity_tradeoff}

Transitions draws on a resource budget whose sustainable rate and ceiling $\mathcal{D}_{\max}$ is set by the substrate. This results in a competition among the three factors.

Transitions draw on a resource budget whose sustainable rate $\mathcal{D}_{\max}$ is set by the substrate and environment. This budget constrains the combinations of workload, toll, and transition flux that the system can sustain.

\begin{proposition}[The Complexity Trade-off]
For homogeneous operation activity in a fixed substrate region, operation within the sustainable budget requires
\begin{equation}
\label{eq:complexity_tradeoff}
    cKf \leq \mathcal{D}_{\max}.
\end{equation}
For $c,K>0$, the corresponding transition-flux bound is
\begin{equation}
\label{eq:update_flux_bound}
    f \leq \frac{\mathcal{D}_{\max}}{cK}.
\end{equation}
At a fixed, binding budget, increasing workload requires a compensating reduction in toll, transition flux, or both.
\end{proposition}

\begin{proof}
Substitute $\mathcal{D}=cKf$ into the budget constraint and divide by $cK>0$.
\end{proof}

For illustration, consider a cell with a fixed ATP budget. If each update requires twice as many ATP-consuming operations while the ATP cost per operation remains unchanged, the same budget supports only half as many updates per unit time. Maintaining the original update rate would require either a lower cost per operation or a larger budget.



\subsection{Operational Action}
\label{subsec:operational_action}

While the operational cost rate describes instantaneous expenditure, evaluating the accumulated burden of distinct histories requires integrating this rate across the system's spatial extent and temporal duration.

Let $\Sigma(t)$ denote the physical or computational elements defining the system boundary at time $t$. Integrating the local cost-rate density over this footprint yields the total instantaneous operational cost:
\begin{equation}
\label{eq:impedance_density}
    Z_{\mathrm{net}}(t)
    = \int_{\Sigma(t)} \mathcal{D}(x,t)\,d\mu(x).
\end{equation}
Here, $\mu$ specifies the substrate measure used to define $\mathcal{D}$. For discrete elements, the counting measure reduces this integral to a sum over the localized elements.

The architectural pillars need not occupy physically disjoint regions of the substrate. Because the six pillars represent distinct functional requirements, their total expenditure partitions exactly into their additive costs (\cref{subsec:cost_of_persistence}), yielding $Z_{\mathrm{net}}(t)=\sum_i Z_i(t)$. This strict accounting ensures that even if a single physical mechanism multiplexes multiple pillars, its operational cost is uniquely apportioned without double-counting.

The \emph{Operational Action} $S_{\Phi}$ of a structural history $\Gamma$ over an observation interval $[0,T]$ is defined as its cumulative operational cost:
\begin{equation}
\begin{split}
\label{eq:operational_action}
    S_{\Phi}[\Gamma;T]
    &= \int_0^T Z_{\mathrm{net}}(t;\Gamma)\,dt \\
    &= \int_0^T\int_{\Sigma(t)} 
      \mathcal{D}(x,t;\Gamma)\,d\mu(x)\,dt.
\end{split}
\end{equation}
Comparing competing systems requires evaluating this action over a common observation interval, a shared resource unit, and a consistent convention for identifying the system boundary.

The dimensions of $S_{\Phi}$ inherit the units of the operational toll $c$. If $c$ represents dissipated energy per operation, $S_{\Phi}$ quantifies macroscopic dissipated energy. If $c$ is defined as a dimensionless action per operation, $S_{\Phi}$ becomes a pure, dimensionless count of state transitions. In arbitrary physical systems, a dimensionless Operational Action can always be recovered by dividing $S_{\Phi}$ by a characteristic reference cost.

While the static architecture identifies the functions instantiated by the system, the Operational Action records the exact thermodynamic burden of their realized operation. This cost accumulates along any history, whether or not the system ultimately persists. Relating this accumulated Operational Action to the probability of survival requires a specific failure model, formalized in \cref{sec:minimum_action}.











\subsection{Entropic Action}
\label{subsec:entropic_action}
Having established the conditions for structural survival, we need a way to compare surviving systems by lifetime. A system's viability cannot be judged pointwise: no single pillar sustains a transition alone, so no single local event determines survival either. Viability is the accumulation across the lifetime. We build that accumulated quantity in two steps: first collapsing the system's spatial extent into one number at each instant, then accumulating that number over its lifetime.

$\mathcal{D}$'s multiplicative, local form and $\bar{Z}_{\mathrm{net}}$'s additive, global one (\cref{subsec:cost_of_persistence}) are reconciled by integrating over the system's spatial footprint $\Sigma(t)$. Because the static architecture partitions the substrate into functionally orthogonal domains corresponding to the six pillars ($\Sigma(t) = \bigcup_i \Sigma_i$), summing the local multiplicative density across these distinct partitions naturally recovers the global additive cost:

\begin{equation}
\begin{split}
\label{eq:impedance_density}
    Z_{\mathrm{net}}(t) &= \int_{\Sigma(t)} \mathcal{D} \, d\mu \\
    &= \sum_i \int_{\Sigma_i} \mathcal{D} \, d\mu \\
    &= \sum_i Z_i
\end{split}
\end{equation}

\noindent with $\mathrm{d}\mu$ the counting measure on the discrete substrate state space, a sum over discrete points, not a continuous integral.

Since $Z_{\mathrm{net}}(t)$ is the total instantaneous dissipation, the cumulative \textbf{Operational Action} is its time integral over the system's lifetime:

\begin{equation}
\label{eq:operational_action}
    S_\Phi = \int_0^{\tau} Z_{\mathrm{net}}(t) \, dt
\end{equation}

\subsubsection{Physical Meaning of the Operational Action}
Because $\mathcal{D}$ counts elementary state transitions per unit substrate measure, $S_\Phi$ is a pure, dimensionless count of irreversible transitions, not yet an energy. When the dimensionless action $S_\Phi$ is projected onto a thermal physical substrate, Landauer's principle supplies the thermodynamic scale $c_0 = k_B T \ln 2$, yielding the physical dissipation analogue $E_{\mathrm{diss}} = c_0 S_\Phi$. On athermal substrates, $S_\Phi$ remains well-defined; only the conversion factor to macroscopic energy changes.

\subsubsection{Necessity of the Architecture for the Action}
\label{subsec:action_requires_architecture}
Minimization of $S_\Phi$ is only possible for systems implementing the complete six-pillar architecture; the Operational Action is not defined independently of it.

\begin{proof}
Persistence requires $\mathcal{D}$ strictly positive and finite, hence $c, K, f > 0$ (\cref{eq:dissipation_density}). This requires all six pillars simultaneously: Toll and Capacity sustain $c > 0$; Map sustains $K > 0$; Protocol, Governor, and Margin jointly sustain $f > 0$. Removing any one pillar collapses the corresponding factor to zero (reproducing the failure modes of \cref{subsec:action_necessity}), so any trajectory that maintains structural coherence necessarily encodes the complete architecture.
\end{proof}

The Operational Action is therefore not defined independently of the architecture.


